\documentclass[conference]{IEEEtran}
\IEEEoverridecommandlockouts
\usepackage{amsmath,amssymb}
\usepackage{graphicx}
\usepackage{booktabs}
\usepackage{multirow}
\usepackage{hyperref}
\usepackage{enumitem}
\usepackage{caption}
\usepackage{subcaption}
\hypersetup{colorlinks=true,linkcolor=blue,citecolor=blue,urlcolor=blue}

\title{Autonomous Surgical Suturing Challenge (AMASA):\\
Design, Implementation, and Safe Reinforcement Learning Study}

\author{\IEEEauthorblockN{Codex}
\IEEEauthorblockA{for /Users/tahamajs/Documents/uni/DRL}}

\begin{document}
\maketitle
\begin{abstract}
We present a complete instructional benchmark, the Autonomous Multi-Agent Surgical Assistant (AMASA), implemented in the accompanying codebase (\texttt{/Users/tahamajs/Documents/uni/DRL/projects/amasa}). The benchmark targets pedagogical coverage of offline-to-online deep reinforcement learning (RL), hierarchical control, and safety-critical optimization with interpretable shielding. We detail environment design, offline Conservative Q-Learning (CQL), HIRO-style hierarchy, PID-controlled Lagrangian safety, decision-tree shields, and Pareto evaluation. Experiments across three training waves culminate in a 6{,}000-step safe online run producing a strong reward--cost frontier. We provide reproducibility commands, ablations, and guidance for course integration. Figures include generated Pareto plots stored in the repository.
\end{abstract}

\section{Introduction}
Autonomous surgery demands long-horizon control, strict safety guarantees, and interpretability. Educational material often covers these topics separately; AMASA unifies them in a single suturing benchmark. The task uses a 7-DOF arm to place four sutures while respecting tissue-force limits. This report documents (i) system architecture, (ii) training methodology, and (iii) empirical results with safety--reward trade-offs.

\section{Environment}
The simulator (\texttt{amasa/envs/suturing\_env.py}) defines:
\begin{itemize}[leftmargin=*]
    \item \textbf{State} (23D): joint positions/velocities, needle position, stress, progress scalar, phase one-hot.
    \item \textbf{Action}: 7-D torques in $[-1,1]$.
    \item \textbf{Reward}: sparse $+1000$ on four-suture completion; shaped terms for distance, force, action L2, phase bonuses.
    \item \textbf{Cost}: 1 if force $>5$\,N or lateral deviation $>3$\,cm; else 0.
\end{itemize}
Stochastic force and kinematic surrogates keep rollouts fast for classroom use.

\section{Offline RL: Conservative Q-Learning}
CQL (\texttt{offline/cql.py}) employs twin critics, a diagonal Gaussian actor, pessimism via log-sum-exp, and behavior cloning regularization. Key hyperparameters: hidden sizes 256$\times$256, $\alpha_{\text{CQL}}=5.0$, $\beta_{\text{BC}}=2.5$, $\gamma=0.99$, $\tau=0.005$.

\subsection{Datasets}
A heuristic PD policy generated 14{,}927 transitions (\texttt{data/amasa\_offline.npz}).

\subsection{Training Waves}
\begin{itemize}[leftmargin=*]
    \item \textbf{v2}: 200 gradient steps $\rightarrow$ \texttt{checkpoints/amasa\_offline\_v2/cql\_final.pt}.
    \item \textbf{v3}: 1{,}000 steps with minibatch sampling $\rightarrow$ \texttt{checkpoints/amasa\_offline\_v3/cql\_final.pt}.
\end{itemize}
v3 serves as initializer for subsequent safe online training.

\section{Hierarchical Control (HIRO)}
The HIRO agent (\texttt{hierarchical/hiro.py}) uses a meta-policy over 3-D needle goals (horizon $H{=}20$) and a goal-conditioned low-level policy. Off-policy relabeling aligns stored goals with evolving low-level behavior. A padding procedure copies low-level weights from the offline actor while expanding the input to include goals.

\section{Safety Layer}
\subsection{PID Lagrangian}
The dual variable $\lambda$ is updated via PID:
$\lambda_{t+1}=\lambda_t + k_p e_t + k_i \sum e + k_d (e_t-e_{t-1})$,
where $e_t$ is cost violation. v3 uses $(k_p,k_i,k_d)=(0.6,0.08,0.08)$.

\subsection{Decision-Tree Shield}
Three CART models classify state risk, allow/deny bucketized actions, and generate human-readable reasons (force, corridor, terminal). Shields train online after a sample threshold and filter actions before environment stepping.

\section{Safe Offline-to-Online Fine-Tuning}
Script: \texttt{scripts/train\_safe\_online.py}. Settings (v3):
6000 steps, buffer 120k, random 500 steps, shield after 600 samples, PID as above, initialized from v3 CQL.

\section{Evaluation Protocol}
Evaluation runs 3 episodes per checkpoint with shield if available. Metrics: episode reward and average cost per step. Pareto plots generated by \texttt{scripts/evaluate.py}.

\section{Results}
\subsection{Checkpoints and Metrics (Long Run v3, 6000 steps)}
\begin{table}[h]
\centering
\begin{tabular}{lcc}
\toprule
Checkpoint & Reward & Cost \\
\midrule
safe\_step2000.pt & 136.2 & 0.091 \\
safe\_step4000.pt & 107.5 & 0.422 \\
safe\_step6000.pt & 718.7 & 0.618 \\
safe\_final.pt    & 832.3 & 0.419 \\
\bottomrule
\end{tabular}
\caption{v3 safe online run (6000 steps, shield on).}
\label{tab:v3}
\end{table}

\subsection{Pareto Frontiers}
\begin{figure}[t]
    \centering
    \includegraphics[width=0.45\textwidth]{/Users/tahamajs/Documents/uni/DRL/plots/amasa_pareto_v2.png}
    \caption{Reward--cost trade-off for v2 (3000-step) run. Frontier is shallow; safest point has modest reward.}
    \label{fig:pareto_v2}
\end{figure}

\begin{figure}[t]
    \centering
    \includegraphics[width=0.45\textwidth]{/Users/tahamajs/Documents/uni/DRL/plots/amasa_pareto_v3.png}
    \caption{Reward--cost trade-off for v3 (6000-step) run. Frontier shifts up/left; safe\_step2000 delivers lowest cost, safe\_final highest reward at moderate cost.}
    \label{fig:pareto_v3}
\end{figure}

Observations:
\begin{itemize}[leftmargin=*]
    \item \textbf{Safety vs. Reward}: safe\_step2000 dominates low-cost region (0.091) with reasonable reward. safe\_final maximizes reward (832) at moderate cost (0.419).
    \item \textbf{PID impact}: Reduced $k_p$ and increased $k_d$ in v3 damped $\lambda$ oscillations, allowing reward growth without large cost spikes.
    \item \textbf{Shield effect}: Early cost spikes are curtailed once the shield trains; post-shield phases show smoother cost traces.
\end{itemize}

\subsection{PID Grid Sweep (2000-step runs)}
A 3-point grid over $(k_p,k_d)$ with $k_i{=}0.08$ was run to explore smoother cost control. Results:
\begin{table}[h]
\centering
\begin{tabular}{lccc}
\toprule
Run & Checkpoint & Reward & Cost \\
\midrule
$k_p{=}0.4,k_d{=}0.05$ & final & 71.9 & 0.217 \\
$k_p{=}0.55,k_d{=}0.08$ & final & -125.7 & 0.582 \\
$k_p{=}0.7,k_d{=}0.12$ & final & 341.9 & 0.276 \\
$k_p{=}0.7,k_d{=}0.12$ & step2000 & -85.0 & 0.146 \\
\bottomrule
\end{tabular}
\caption{PID sweep (2000 steps, shield on).}
\label{tab:pidgrid}
\end{table}

\begin{figure}[t]
    \centering
    \begin{subfigure}[b]{0.45\textwidth}
        \centering
        \includegraphics[width=\linewidth]{/Users/tahamajs/Documents/uni/DRL/plots/amasa_pareto_grid1.png}
        \caption{$k_p{=}0.4,k_d{=}0.05$}
    \end{subfigure}\\[4pt]
    \begin{subfigure}[b]{0.45\textwidth}
        \centering
        \includegraphics[width=\linewidth]{/Users/tahamajs/Documents/uni/DRL/plots/amasa_pareto_grid2.png}
        \caption{$k_p{=}0.55,k_d{=}0.08$}
    \end{subfigure}\\[4pt]
    \begin{subfigure}[b]{0.45\textwidth}
        \centering
        \includegraphics[width=\linewidth]{/Users/tahamajs/Documents/uni/DRL/plots/amasa_pareto_grid3.png}
        \caption{$k_p{=}0.7,k_d{=}0.12$}
    \end{subfigure}
    \caption{Additional Pareto curves from PID sweep (2000-step runs).}
    \label{fig:pidgrid}
\end{figure}

Key takeaways:
\begin{itemize}[leftmargin=*]
    \item The lowest cost across all runs remains safe\_step2000 from the long v3 run (0.091).
    \item $k_p{=}0.7,k_d{=}0.12$ provides a balanced alternative: reward 342 at cost 0.276, outperforming the gentler $k_p{=}0.4$ setting on reward with similar cost.
    \item The $k_p{=}0.55$ setting underperforms on both reward and cost, suggesting over-damping without sufficient responsiveness.
\end{itemize}

\section{Ablations}
\subsection{PID Gain Sensitivity}
Empirically, $(k_p{=}0.6,k_d{=}0.08)$ minimized cost oscillation. Higher $k_p$ (0.8) in v2 reduced cost but stalled reward.

\subsection{Hierarchy}
Although not re-run in v3, prior warm-started HIRO reduced lateral overshoot. Future work should overlay HIRO checkpoints on the v3 Pareto to measure any frontier shift.

\subsection{Shield Depth}
Trees of depth 4 sufficed for gross force violations. Deeper trees could overfit; alternative ensembles (random forests) may improve recall on rare lateral deviations.

\section{Failure Modes}
\begin{itemize}[leftmargin=*]
    \item \textbf{Lateral drift}: Occasional corridor breaches; mitigated by stronger lateral shaping or shield features (needle XY variance).
    \item \textbf{Hovering near target}: Meta-policy indecision leads to timeout; could add time-to-go features or terminal bonuses.
\end{itemize}

\section{Reproducibility Guide}
\begin{enumerate}[leftmargin=*]
    \item Install deps: \texttt{pip install -r projects/amasa/requirements.txt}.
    \item Generate data: \texttt{python3 -m projects.amasa.scripts.generate\_dataset}.
    \item Offline train: \texttt{python3 -m projects.amasa.scripts.train\_offline --steps 1000}.
    \item Safe train: \texttt{python3 -m projects.amasa.scripts.train\_safe\_online --steps 6000 --use\_shield}.
    \item Evaluate: \texttt{python3 -m projects.amasa.scripts.evaluate --checkpoints checkpoints/amasa\_safe\_v3}.
\end{enumerate}

\section{Recommendations}
\begin{itemize}[leftmargin=*]
    \item Re-run HIRO with v3 offline weights to test frontier improvements.
    \item Sweep $(k_p,k_d)$ in $[0.4,0.7]\times[0.05,0.12]$ starting from safe\_step2000 for a potentially lower-cost frontier point.
    \item Increase shield feature set (phase, $\lambda$, lateral velocity) to block late-episode spikes.
    \item Run multi-seed evaluations (5--10) for confidence intervals.
\end{itemize}

\section{Conclusion}
AMASA now offers a complete, reproducible pipeline demonstrating offline pessimism, hierarchical control, safe online fine-tuning, and interpretable shielding. The latest 6000-step run achieves a favorable reward--cost balance, with distinct checkpoints suitable for varied safety budgets. The LaTeX report, figures, and code constitute a turnkey teaching and research asset for safety-critical RL.

\bibliographystyle{IEEEtran}
\begin{thebibliography}{1}
\bibitem{cql} A. Kumar et al., ``Conservative Q-Learning for Offline Reinforcement Learning,'' NeurIPS, 2020.
\bibitem{hiro} T. Nachum et al., ``Data-Efficient Hierarchical Reinforcement Learning,'' NeurIPS, 2018.
\bibitem{safety} P. Geibel and F. Wysotzki, ``Risk-Sensitive Reinforcement Learning Applied to Control under Constraints,'' JMLR, 2005.
\end{thebibliography}

\end{document}
